% ============================================================
% House style: copied from
%   ../../PCORTC_derivative/20260920_karuizawa_seminar_slide/beamer_header.tex
% with the header colours left at the moloch defaults and without LuaTeX-ja.
% ============================================================
\usetheme[numbering=fraction,sectionpage=progressbar, subsectionpage=progressbar]{moloch}
\molochset{block=fill}
\setbeamercovered{transparent}
\setbeamertemplate{items}[square]
\definecolor{chiba}{RGB}{96,39,40}
\definecolor{chibasub}{RGB}{199,0,43}
% Header colours (frame title, title separator, progress bar): moloch defaults.
\setbeamercolor{block title}{
  % use=palette primary,
  % parent=palette primary,
  fg = black,
  bg = chibasub!20
}
\setbeamercolor{block body}{
bg = block title.bg!20!normal text.bg
}%
\setbeamercolor{block title alerted}{
  fg = black,
  bg = chibasub!50
}
\setbeamercolor{block body alerted}{
bg = block title alerted.bg!20!normal text.bg
}
\setbeamercolor{block title example}{
  % use=palette primary,
  % parent=palette primary,
  fg = example text.fg!20!black,
  bg = example text.fg!20!chibasub!30
}
\setbeamercolor{block body example}{
  bg = block title example.bg!20!normal text.bg
}
\addtobeamertemplate{proof begin}{%
  \setbeamercolor{block title}{
    use=palette primary,
    parent=palette primary,
    fg = black,
    bg = chibasub!50
  }%
  \setbeamercolor{block body}{
    bg = block title.bg!20!normal text.bg
  }%
}{}%
\AtBeginEnvironment{theorem}{%
  \setbeamercolor{block title}{
    fg = black, bg = chibasub!50
  }%
  \setbeamercolor{block body}{
    bg = block title.bg!20!normal text.bg
  }%
}
\AtBeginEnvironment{definition}{%
  \setbeamercolor{block title}{
    fg = black, bg = chibasub!50
  }%
  \setbeamercolor{block body}{
    bg = block title.bg!20!normal text.bg
  }%
}
\AtBeginEnvironment{lemma}{%
  \setbeamercolor{block title}{
    fg = black, bg = chibasub!30
  }%
  \setbeamercolor{block body}{
    bg = block title.bg!20!normal text.bg
  }%
}
\AtBeginEnvironment{corollary}{%
  \setbeamercolor{block title}{
    fg = black, bg = chibasub!30
  }%
  \setbeamercolor{block body}{
    bg = block title.bg!20!normal text.bg
  }%
}
\AtBeginEnvironment{proposition}{%
  \setbeamercolor{block title}{
    fg = black, bg = chibasub!30
  }%
  \setbeamercolor{block body}{
    bg = block title.bg!20!normal text.bg
  }%
}
\AtBeginEnvironment{problem}{%
  \setbeamercolor{block title}{
    fg = black, bg = chibasub!30!blue!20  % light enough for gray citations in the title
  }%
  \setbeamercolor{block body}{
    bg = block title.bg!20!normal text.bg
  }%
}

\usefonttheme[onlymath]{serif}

% begin bibliography icons
\setbeamertemplate{bibliography item}{%
  \ifboolexpr{ test {\ifentrytype{book}} or test {\ifentrytype{mvbook}}
    or test {\ifentrytype{collection}} or test {\ifentrytype{mvcollection}}
    or test {\ifentrytype{reference}} or test {\ifentrytype{mvreference}} }
    {\setbeamertemplate{bibliography item}[book]}
    {\ifentrytype{online}
       {\setbeamertemplate{bibliography item}[online]}
       {\setbeamertemplate{bibliography item}[article]}}%
  \usebeamertemplate{bibliography item}}
% end

% begin enum icon
\usepackage{enumitem}
\newcommand{\labelitemi}{\textbullet}
\newcommand{\labelitemii}{\textbullet}
\newcommand{\labelenumi}{\arabic{enumi}.}
% end

% 吹き出し
\newcommand{\tikzlabel}[1]{\resizebox{0pt}{!}{\begin{tikzpicture}[remember picture] \node(#1){}; \end{tikzpicture}}}
\tikzstyle{mycallout}=[rectangle callout, white, fill=red!80, rounded corners]
\tikzstyle{mycallouty}=[rectangle callout, black, fill=yellow!80, rounded corners]
\tikzstyle{mycalloutnote}=[rectangle callout, black, draw=black, fill=lightgray!40, rounded corners]
\tikzstyle{mycalloutnote2}=[rectangle, black, draw=black, fill=lightgray!40, rounded corners]
%強調
\usepackage[most]{tcolorbox}
\tcbset{on line,
    boxsep=2pt, left=0pt,right=0pt,top=0pt,bottom=0pt,
    colframe=red!80,colback=red!80,
    highlight math style={enhanced}
}
\newcommand{\superemph}[1]{\tcbox{\textcolor{white}{#1}}}
%highlight
\newcommand<>\hlbox[2]{\alt#3{\makebox[\dimexpr\width-2\fboxsep]{\colorbox{#1}{#2}}}{#2}}

\makeatletter
\setlength{\moloch@progressonsectionpage@linewidth}{2.pt}
\makeatother
\AtBeginSection[]
{
    \begin{frame}[plain, c, noframenumbering]
        \frametitle{\faListUl \quad Outline}
        \usebeamercolor[fg]{section title}
        % \usebeamerfont{section title}
        % \insertsectionhead\\[-1ex]
        \tableofcontents[
            currentsection,
            sectionstyle=show/shaded,
            subsectionstyle=show/shaded/hide,
            subsubsectionstyle=show/shaded/hide
        ]
        \usebeamertemplate*{progress bar in section page}
    \end{frame}
}
\AtBeginSubsection[]
{
    % \setbeamercolor{background canvas}{bg=white}
    \begin{frame}[plain, c, noframenumbering]
        \frametitle{\faListUl \quad Outline}
        \tableofcontents[
            currentsection,
            sectionstyle=show/shaded,
            subsectionstyle=show/shaded/hide,
            subsubsectionstyle=show/shaded/hide
        ]
        \usebeamertemplate*{progress bar in section page}
    \end{frame}
}

\AtBeginSubsubsection[]
{
    % \setbeamercolor{background canvas}{bg=white}
    \begin{frame}<handout:0>[plain, c, noframenumbering]
        \frametitle{\faListUl \quad Outline}
        \tableofcontents[
            currentsection,
            sectionstyle=show/shaded,
            subsectionstyle=show/shaded/hide,
            subsubsectionstyle=show/shaded/hide
        ]
        \usebeamertemplate*{progress bar in section page}
    \end{frame}
}


% align* 
\AtBeginDocument{
    \abovedisplayskip     =0.5\abovedisplayskip
    \abovedisplayshortskip=0.5\abovedisplayshortskip
    \belowdisplayskip     =0.5\belowdisplayskip
    \belowdisplayshortskip=0.5\belowdisplayshortskip}


\usepackage{fontspec, fontawesome5}

\newcommand{\myIconCoffeeBreak}{\textcolor{green}{\faIcon{mug-hot}}}
\newcommand{\myIconCaution}{\textcolor{yellow}{\faIcon{exclamation-triangle}}}
\newcommand{\myIconTask}{\textcolor{cyan}{\faIcon{list}}}
\newcommand{\myIconVertical}{\raisebox{-.15ex}{\rotatebox{90}{\faIcon{minus}}}}
\newcommand{\myIconQuestion}{\textcolor{cyan}{\scalebox{1}[1]{\faIcon{question}}}}
\newcommand{\myIconQuestionLine}{\textcolor{cyan}{\myIconVertical}}


\def\proofname{\translate{Proof}}
\newtheorem{proposition}{\translate{Proposition}}
\newtheorem{notation}{\translate{Notation}}
\newtheorem{task}{\translate{Task}}
\theoremstyle{example}
\newtheorem{exampleproblem}{\translate{Example Problem}}
\theoremstyle{remark}
\newtheorem{remark}{\translate{Remark}}
\newtheorem{conjecture}{\translate{Conjecture}}

\uselanguage{english}
\languagepath{english}
\deftranslation[to=japanese]{Theorem}{\textcolor{red}{\myIconVertical}定理}
\deftranslation[to=english]{Theorem}{\textcolor{red}{\myIconVertical}Thm.}
\deftranslation[to=japanese]{Lemma}{\textcolor{red!50}{\myIconVertical}補題}
\deftranslation[to=english]{Lemma}{\textcolor{red!50}{\myIconVertical}Lem.}
\deftranslation[to=japanese]{Proposition}{\textcolor{red!30}{\myIconVertical}命題}
\deftranslation[to=english]{Proposition}{\textcolor{red!30}{\myIconVertical}Prop.}
\deftranslation[to=japanese]{Corollary}{\textcolor{red!30}{\myIconVertical}系}
\deftranslation[to=english]{Corollary}{\textcolor{red!30}{\myIconVertical}Cor.}
\deftranslation[to=japanese]{Proof}{証明}
\deftranslation[to=english]{Proof}{Proof}

\deftranslation[to=japanese]{Definition}{\textcolor{black}{\myIconVertical}定義}
\deftranslation[to=english]{Definition}{\textcolor{black}{\myIconVertical}Def.}
\deftranslation[to=english]{Fact}{\textcolor{black}{\myIconVertical}Fact}
\deftranslation[to=japanese]{Definitions}{\textcolor{black}{\myIconVertical}定義}
\deftranslation[to=japanese]{Notation}{表記}
\deftranslation[to=japanese]{Example}{例}
\deftranslation[to=japanese]{Examples}{例}
\deftranslation[to=japanese]{Problem}{\myIconQuestion 問題}
\deftranslation[to=japanese]{Note}{\myIconCaution ノート}
\deftranslation[to=japanese]{Remark}{\myIconCaution 注記}
% \AtBeginEnvironment{task}{%
%     \setbeamercolor{block body}{bg=normal text.bg!80!fg}
%     \setbeamercolor{block title}{bg=red!80,fg=black!2}
% }
\deftranslation[to=japanese]{Task}{\myIconTask 課題}
\deftranslation[to=japanese]{Example Problem}{\myIconQuestion 例題}
\deftranslation[to=japanese]{Solution}{解}
\deftranslation[to=japanese]{Fact}{事実}

\usepackage[backend=biber, uniquelist=false, maxcitenames = 2, maxbibnames=8, style=alphabetic, maxalphanames=3, minalphanames=3, sorting = none, firstinits=true, sortcites=false]{biblatex}
\DeclareNameAlias{sortname}{given-family}
\DeclareNameAlias{default}{given-family}

% https://tex.stackexchange.com/questions/73136/make-specific-author-bold-using-biblatex
\usepackage{xpatch}
\makeatletter
\newbibmacro*{name:bold}[2]{%
  \edef\blx@tmp@name{\expandonce#1, \expandonce#2}%
  \def\do##1{\ifdefstring{\blx@tmp@name}{##1}{\bfseries\listbreak}{}}%
  \dolistloop{\boldnames}}
\newcommand*{\boldnames}{}
\makeatother
\xpretobibmacro{name:family}{\begingroup\usebibmacro{name:bold}{#1}{#2}}{}{}
\xpretobibmacro{name:given-family}{\begingroup\usebibmacro{name:bold}{#1}{#2}}{}{}
\xpretobibmacro{name:family-given}{\begingroup\usebibmacro{name:bold}{#1}{#2}}{}{}
\xpretobibmacro{name:delim}{\begingroup\normalfont}{}{}

\xapptobibmacro{name:family}{\endgroup}{}{}
\xapptobibmacro{name:given-family}{\endgroup}{}{}
\xapptobibmacro{name:family-given}{\endgroup}{}{}
\xapptobibmacro{name:delim}{\endgroup}{}{}

% \usepackage{makeidx}
% \makeindex
% \usepackage[xcolor, hyperref, notion, makeidx]{knowledge}
% ? -> ../header.tex


\newcommand{\xmark}{%
    \tikz[scale=0.23] {
        \draw[line width=0.7,line cap=round] (0,0) to [bend left=6] (1,1);
        \draw[line width=0.7,line cap=round] (0.2,0.95) to [bend right=3] (0.8,0.05);
    }}
\newcommand{\cmark}{%
    \tikz[scale=0.23] {
        \draw[line width=0.7,line cap=round] (0.25,0) to [bend left=10] (1,1);
        \draw[line width=0.8,line cap=round] (0,0.35) to [bend right=1] (0.23,0);
    }}

\newcommand*\circled[1]{\tikz[baseline=(char.base)]{
        \node[shape=circle,draw,inner sep=1pt] (char) {#1};}}


\newcommand{\defbox}[1]{\colorbox{blue!30}{#1}}
\newcommand{\EMPH}[1]{\emph{\textcolor{orange}{#1}}}

\usepackage{makecell}
\usepackage{colortbl}

\usepackage{annotate-equations}